\setcounter{section}{8}

  \zwischenueberschrift{Soal latihan}

\inputaufgabe
{}
{

Misalkan
\mathkor {} {M} {dan} {N} {}
adalah
\definitionsverweis {manifold diferensiabel}{}{}
dan
\maabbdisp {\varphi} { M } { N
} {}
sebuah pemetaan. Misalkan
\mathl{M= \bigcup_{i \in I} U_i}{} suatu
\definitionsverweis {tutupan terbuka}{}{}
dari $M$. Tunjukkan bahwa $\varphi$
\definitionsverweis {diferensiabel}{}{}
jika dan hanya jika semua pembatasan
\mathl{\varphi_i= \varphi \vert_{U_i }}{} diferensiabel.

}
{} {}

\inputaufgabe
{}
{

Misalkan $M$ sebuah
\definitionsverweis {manifold diferensiabel}{}{}
dan
\maabbdisp {\alpha} { U } { V
} {}
sebuah peta
\zusatzklammer {jadi
\mathl{U \subseteq M}{} dan
\mathl{V \subseteq \R^n}{} terbuka} {} {.}
Tunjukkan bahwa $\alpha$ adalah sebuah
\definitionsverweis {difeomorfisme}{}{}.

}
{} {}

\inputaufgabe
{}
{

Misalkan $M$ sebuah
\definitionsverweis {manifold diferensiabel}{}{} dan
\maabbdisp {f,g} { M } {\R
} {}
\definitionsverweis {fungsi-fungsi diferensiabel}{}{} pada $M$. Buktikan pernyataan-pernyataan berikut. \aufzaehlungvier{Pemetaan
\maabbeledisp {f \times g} { M } {\R^2
} { x } {(f(x),g(x))
} {,} diferensiabel.
}{ $f+g$
\definitionsverweis {diferensiabel}{}{.}
}{ $f \cdot g$
\definitionsverweis {diferensiabel}{}{.}
}{Jika $f$ tidak memiliki titik nol, maka
\mathl{f^{-1}}{} juga diferensiabel.
}

}
{} {}

\inputaufgabe
{}
{

Misalkan
\mathl{S^2 \subseteq \R^3}{} adalah
\definitionsverweis {sfera}{}{.}
Dengan menggunakan
\definitionsverweis {peta-peta stereografis}{}{,}
tunjukkan bahwa pembatasan koordinat ruang
\mathl{x,y,z}{} pada sfera merupakan
\definitionsverweis {fungsi-fungsi diferensiabel}{}{}.

}
{} {}

\inputaufgabe
{}
{

Misalkan
\mathkor {} {M} {dan} {N} {}
adalah
\definitionsverweis {manifold diferensiabel}{}{}
dan
\maabbdisp {\varphi} { M } { N
} {}
sebuah
\definitionsverweis {pemetaan diferensiabel}{}{.}
Tunjukkan bahwa pemetaan ini menginduksi sebuah
\definitionsverweis {homomorfisme gelanggang}{}{}
\maabbeledisp {\varphi^*} {C^1(N,\R)} {C^1(M,\R)
} { f } {f \circ \varphi
} {,}
melalui penarikan balik.

}
{} {}

\inputaufgabe
{}
{

Tunjukkan bahwa untuk
\mathl{m \leq n}{}, penyematan subruang
\mathl{\R^m}{} ke dalam $\R^n$ yang diberikan oleh
\mathl{(x_1 , \ldots , x_m) \mapsto (x_1 , \ldots , x_m,0 , \ldots , 0)}{}
\definitionsverweis {diferensiabel hingga orde berapa pun}{}{}.

}
{} {}

\inputaufgabe
{}
{

Tunjukkan bahwa permukaan silinder terbuka
\mathl{S^1 \times {]0,1[}}{},
\mathl{S^1 \times \R}{,} bidang yang ditusuk
\mathl{\R^2 \setminus \{(0,0)\}}{}, dan
\mathl{S^2 \setminus \{N,S\}}{}
saling \definitionsverweis {difeomorfik}{}{}.

}
{} {}

Soal berikut menggunakan definisi di bawah ini.

Sebuah fungsi
\maabbdisp {f} { {\mathbb K}^n } { {\mathbb K}
} {}
disebut
\definitionswort {homogen berderajat}{}
$d$, jika untuk setiap titik
\mathl{P \in {\mathbb K}^n}{} dan setiap
\mathl{\lambda \in {\mathbb K}}{} berlaku hubungan

\mavergleichskettedisp
{\vergleichskette
{f(\lambda P)
}
{ =} { \lambda^d f(P)
}
{ } {
}
{ } {
}
{ } {
}
}
{}{}{}
untuk semua pilihan tersebut.

\inputaufgabe
{}
{

Misalkan
\maabbdisp {f} {\R^n } {\R
} {}
sebuah
\definitionsverweis {fungsi homogen}{}{} yang
\definitionsverweis {terdiferensialkan secara kontinu}{}{,}
dan yang pada
\definitionsverweis {serat}{}{}
$F$ di atas
\mathl{a \neq 0}{}
bersifat \definitionsverweis {reguler}{}{.}
Tunjukkan bahwa setiap serat di atas
\mathl{b \neq 0,\ b=\lambda^d a}{}
untuk suatu skalar real tak nol,
\definitionsverweis {difeomorfik}{}{} dengan $F$ dan merupakan sebuah
\definitionsverweis {manifold}{}{}.

Catatan edisi: Syarat penskalaan real ditambahkan. Klaim sumber untuk semua nilai tak nol gagal pada derajat genap ketika tanda nilai serat berbeda; misalnya, fungsi kuadrat norma mempunyai serat positif berbentuk bola tetapi serat negatif kosong.

}
{} {}

\inputaufgabe
{}
{

Misalkan
\mathl{]a,b[}{} sebuah
\definitionsverweis {interval terbuka}{}{}
dan
\maabbdisp {f} {]a,b[} {\R_{+}
} {}
sebuah
\definitionsverweis {fungsi yang terdiferensialkan secara kontinu}{}{.}
Misalkan $M$ adalah permukaan benda putar terkait. Tunjukkan bahwa himpunan ini merupakan manifold yang difeomorfik dengan silinder terbuka.

}
{} {}

\inputaufgabe
{}
{

Tunjukkan bahwa permukaan elipsoid dan permukaan bola satuan
$C^\infty$-\definitionsverweis {difeomorfik}{}{}.

}
{} {}

\inputaufgabegibtloesung
{}
{

Berikan sebuah contoh
\definitionsverweis {manifold diferensiabel}{}{}
\definitionsverweis {berdimensi dua}{}{} yang
\definitionsverweis {terhubung}{}{},
sebutlah $M$, dan sebuah titik

\mavergleichskette
{\vergleichskette
{ P
}
{ \in }{ M
}
{  }{
}
{    }{
}
{   }{
}
}
{}{}{}
sedemikian sehingga
\mathkor {} {M} {dan} {M \setminus \{P\}} {}
\definitionsverweis {difeomorfik}{}{}
satu sama lain.

}
{} {}

\inputaufgabe
{}
{

Tunjukkan bahwa
\definitionsverweis {pemetaan antipodal}{}{}
\maabbeledisp {} {S^n } { S^n
} {(x_1 , \ldots , x_{n+1})} {(-x_1 , \ldots , -x_{n+1})
} {,}
adalah sebuah
\definitionsverweis {difeomorfisme}{}{} \definitionsverweis {tanpa titik tetap}{}{}
yang inversnya adalah dirinya sendiri.

}
{} {}

\inputaufgabegibtloesung
{}
{

Tunjukkan bahwa sebuah
\definitionsverweis {manifold diferensiabel}{}{}
$M$ berdimensi

\mavergleichskette
{\vergleichskette
{ n
}
{ \geq }{ 1
}
{  }{
}
{    }{
}
{   }{
}
}
{}{}{}
memiliki tak hingga banyak
\definitionsverweis {difeomorfisme}{}{}
\maabbdisp {\varphi} { M } { M
} {}
yang berbeda.

}
{} {}

Soal-soal berikut dimaksudkan untuk menjelaskan mengapa manifold dirumuskan dengan tutupan terbuka.

Misalkan
\mathl{\left(  x_n  \right)_{n \in  \N  }}{} sebuah
\definitionsverweis {barisan}{}{}
dalam suatu
\definitionsverweis {ruang topologis}{}{}
$X$. Kita mengatakan bahwa, untuk suatu

\mavergleichskette
{\vergleichskette
{ x
}
{ \in }{ X
}
{  }{
}
{    }{
}
{   }{
}
}
{}{}{}
barisan tersebut \definitionswort {konvergen ke titik itu}{,} jika sifat berikut terpenuhi.

Untuk setiap
\definitionsverweis {lingkungan terbuka}{}{}

\mavergleichskette
{\vergleichskette
{ U
}
{ \subseteq }{ X
}
{  }{
}
{    }{
}
{   }{
}
}
{}{}{}
yang memuat $x$, terdapat suatu

\mavergleichskette
{\vergleichskette
{ n_0
}
{ \in }{ \N
}
{  }{
}
{    }{
}
{   }{
}
}
{}{}{}
sedemikian sehingga untuk semua

\mavergleichskette
{\vergleichskette
{ n
}
{ \geq }{ n_0
}
{  }{
}
{    }{
}
{   }{
}
}
{}{}{}
suku barisan $x_n$ termasuk dalam $U$.

Dalam hal ini $x$ disebut \definitionswort {nilai batas}{} atau \definitionswort {limit}{} barisan tersebut. Kita juga menuliskannya sebagai

\mavergleichskettedisp
{\vergleichskette
{  \lim_{n \rightarrow \infty}  x_n
}
{ =} { x
}
{ } {
}
{ } {
}
{ } {
}
}
{}{}{.}
Jika barisan memiliki limit, kita juga mengatakan bahwa barisan itu \definitionswort {konvergen}{}
\zusatzklammer {tanpa menyebut nilai limit tertentu} {} {,}
dan jika tidak, bahwa barisan itu \definitionswort {divergen}{.}

\inputaufgabe
{}
{

Misalkan $M$ sebuah
\definitionsverweis {manifold topologis}{}{}
dan
\mathl{\left(   x_n  \right)_{n \in \N }}{} sebuah
\definitionsverweis {barisan}{}{}
dalam $M$. Tunjukkan bahwa barisan tersebut
\definitionsverweis {konvergen}{}{,}
jika dan hanya jika terdapat sebuah
\definitionsverweis {domain peta}{}{}
dari $M$ yang memuat hampir semua suku barisan dan sedemikian sehingga barisan citra yang bersesuaian konvergen di dalam citra peta.

}
{} {}

\inputaufgabe
{}
{

Misalkan $M$ sebuah
\definitionsverweis {manifold topologis}{}{}
dan
\maabbdisp {\alpha_i} {U_i } { V_i
} {}
sebuah keluarga
\definitionsverweis {peta}{}{}
dengan
\definitionsverweis {pemetaan transisi}{}{}
\maabbdisp {\varphi_{ij} = \alpha_j \circ (\alpha_i)^{-1}} {V_i \cap \alpha_i(U_i \cap U_j) } { V_j \cap \alpha_j(U_i \cap U_j)
} {.}
Tunjukkan bahwa manifold $M$ dapat direkonstruksi dari keluarga

\mathbed {V_i} {}
 {i \in I} {}
 {} {} {} {,}
himpunan-himpunan bagian
\mathl{V_{ij} \subseteq V_i}{} dan pemetaan transisi
\maabbdisp {\varphi_{ij}} {V_{ij}} { V_{ji}
} {}
yang bersesuaian.

a) Pada

\mavergleichskettedisp
{\vergleichskette
{N
}
{ \defeq} { \biguplus_{i \in I} V_i
}
{ } {
}
{ } {
}
{ } {
}
}
{}{}{}
perhatikan
\definitionsverweis {relasi ekuivalensi}{}{,}
yang menyamakan dua titik
\mathl{P\in V_i}{} dan
\mathl{Q \in V_j}{} jika
\mathl{\varphi_{ij}}{} memetakan titik pertama ke titik kedua.

b) Lengkapi himpunan hasil bagi
\mathl{N/ \sim}{} dengan sebuah topologi yang sesuai.

c) Definisikan peta-peta pada
\mathl{N/ \sim}{}.

d) Tunjukkan bahwa
\mathkor {} {M} {dan} {N/ \sim} {}
\definitionsverweis {homeomorfik}{}{}.

}
{} {}

Prosedur konstruksi manifold yang dijelaskan dalam soal sebelumnya berlaku untuk suatu keluarga himpunan bagian terbuka di dalam $\R^n$ dengan pemetaan transisi yang memenuhi syarat koklus dari
Soal 7.11.
Namun, ruang topologis yang dihasilkan tidak serta-merta merupakan ruang Hausdorff. Konstruksi ini disebut \stichwort {perekatan terbuka} {} ruang-ruang.

\inputaufgabe
{}
{

Kita mengambil dua salinan garis riil, yaitu

\mavergleichskette
{\vergleichskette
{G_1
}
{ = }{\R
}
{  }{
}
{    }{
}
{   }{
}
}
{}{}{}
dan

\mavergleichskette
{\vergleichskette
{G_2
}
{ = }{\R
}
{  }{
}
{    }{
}
{   }{
}
}
{}{}{}
beserta pemetaan perekatan
\maabbeledisp {\varphi} {G_1 \setminus \{0\} } { G_2 \setminus \{0\}
} { x } { x^{-1}
} {.}
Misalkan $M$ ruang topologis yang dihasilkan menurut konstruksi yang dijelaskan dalam
Soal 8.15.
Tunjukkan bahwa $M$
\definitionsverweis {homeomorfik}{}{}
dengan
$1$-\definitionsverweis {sfera}{}{}.

}
{} {}

\inputaufgabe
{}
{

Kita mengambil dua salinan garis riil, yaitu

\mavergleichskette
{\vergleichskette
{G_1
}
{ = }{\R
}
{  }{
}
{    }{
}
{   }{
}
}
{}{}{}
dan

\mavergleichskette
{\vergleichskette
{G_2
}
{ = }{\R
}
{  }{
}
{    }{
}
{   }{
}
}
{}{}{}
beserta pemetaan perekatan
\maabbdisp {\operatorname{Id}} {G_1 \setminus \{0\} } { G_2 \setminus \{0\}
} {.}
Misalkan $M$ ruang topologis yang dihasilkan menurut konstruksi yang dijelaskan dalam
Soal 8.15.
Tunjukkan bahwa $M$ bukan
\definitionsverweis {manifold}{}{}.

}
{} {}

  \zwischenueberschrift{Soal untuk dikumpulkan}

Dalam soal berikut, tafsirkan ${\mathbb C}$ sebagai $\R^2$.

\inputaufgabe
{4}
{

Perhatikan pemetaan
\maabbeledisp {} { {\mathbb C}^2} { {\mathbb C}
} {(z,w)} { zw
} {.}
Untuk titik manakah
\mathl{u \in {\mathbb C}}{} serat di atas $u$ merupakan manifold? Dalam setiap kasus, berikan deskripsi kelas difeomorfisme yang sesederhana mungkin.

}
{} {}

\inputaufgabe
{6}
{

Misalkan diberikan dua titik
\mathkor {} {P} {dan} {Q} {}
pada
\definitionsverweis {permukaan bola satuan}{}{.}
Tunjukkan bahwa terdapat sebuah
\definitionsverweis {difeomorfisme}{}{}
dari permukaan bola itu ke dirinya sendiri yang memetakan $P$ ke $Q$.

}
{} {}

\inputaufgabe
{4 (1+1+2)}
{

Misalkan $M$ sebuah
\definitionsverweis {manifold diferensiabel}{}{.}
Untuk setiap himpunan bagian terbuka

\mavergleichskette
{\vergleichskette
{ U
}
{ \subseteq }{ M
}
{  }{
}
{    }{
}
{   }{
}
}
{}{}{}
kita perhatikan himpunan
\mathl{C^1(U,\R)}{} yang terdiri atas
\definitionsverweis {fungsi-fungsi diferensiabel}{}{}
pada $U$. Misalkan

\mavergleichskette
{\vergleichskette
{ M
}
{ = }{ \bigcup_{i \in I} U_i
}
{  }{
}
{    }{
}
{   }{
}
}
{}{}{}
sebuah
\definitionsverweis {tutupan terbuka}{}{.}
\aufzaehlungdreiabc{Tunjukkan bahwa jika

\mavergleichskette
{\vergleichskette
{ V
}
{ \subseteq }{ U
}
{  }{
}
{    }{
}
{   }{
}
}
{}{}{}
terbuka dan

\mavergleichskette
{\vergleichskette
{ f
}
{ \in }{ C^1(U,\R)
}
{  }{
}
{    }{
}
{   }{
}
}
{}{}{}
maka
\definitionsverweis {pembatasan}{}{}
$f \vert_V$ juga termasuk dalam
\mathl{C^1(V,\R)}{}.
}{Misalkan

\mavergleichskette
{\vergleichskette
{ f
}
{ \in }{ C^1(M,\R)
}
{  }{
}
{    }{
}
{   }{
}
}
{}{}{.}
Tunjukkan bahwa

\mavergleichskette
{\vergleichskette
{ f
}
{ = }{ 0
}
{  }{
}
{    }{
}
{   }{
}
}
{}{}{}
jika dan hanya jika semua pembatasan

\mavergleichskette
{\vergleichskette
{ f \vert_{U_i }
}
{ = }{ 0
}
{  }{
}
{    }{
}
{   }{
}
}
{}{}{}
berlaku.
}{Misalkan diberikan suatu keluarga

\mavergleichskette
{\vergleichskette
{ f_i
}
{ \in }{ C^1(U_i,\R)
}
{  }{
}
{    }{
}
{   }{
}
}
{}{}{}
fungsi-fungsi yang memenuhi \anfuehrung{syarat kecocokan}{}

\mavergleichskettedisp
{\vergleichskette
{ f_i \vert_{U_i \cap U_j }
}
{ =} { f_j \vert_{U_i \cap U_j }
}
{ } {
}
{ } {
}
{ } {
}
}
{}{}{}
untuk semua $i,j$. Tunjukkan bahwa terdapat suatu

\mavergleichskette
{\vergleichskette
{ f
}
{ \in }{ C^1(M,\R)
}
{  }{
}
{    }{
}
{   }{
}
}
{}{}{}
yang memenuhi

\mavergleichskette
{\vergleichskette
{ f \vert_{U_i }
}
{ = }{ f_i
}
{  }{
}
{    }{
}
{   }{
}
}
{}{}{}
untuk semua $i$.
}

}
{} {}

\inputaufgabe
{5}
{

Misalkan $M$ sebuah
\definitionsverweis {manifold diferensiabel}{}{,} yang memiliki setidaknya dua titik. Tunjukkan bahwa terdapat
\definitionsverweis {fungsi-fungsi diferensiabel}{}{}
\maabbdisp {f,g} { M } {\R
} {}
dengan
\mathl{f,g \neq 0}{,} tetapi
\mathl{fg=0}{.}

}
{} {}
